\section{失守之后：恢复、接收与记忆的不同速度}

\subsection{城池易手后的第一年}

一座城被攻下或接受新政权，史书通常以日期、诏令和官员任命为其结局。对城中与周边居民而言，第一年更可能由粮食、治安、田地、亲属和谁来征税构成。新军需要驻扎和补给，旧有官吏与地方中介要决定去留，商人要判断道路是否安全，逃离者要寻找落脚点。政治的转换有明确时点，生活秩序的转换却需要更长时间。

宋方材料在失守后常使用忠、降、义等词汇，元方材料则倾向强调安抚、任命和秩序恢复。两类文字都保存了重要信息，也都带有要说服读者的目的。忠义语言不能告诉我们所有守城者的动机，安抚文书也不能证明每一位居民接受新秩序。将它们放回粮道、城门、仓储和亲属网络，才能避免让政治分类代替地方经验。

\subsection{重建不是回到原处}

战后修城、修堤、重开市场或恢复寺院，常被称为“恢复”。这个词容易让人以为先前的秩序能够原样回来。实际的修复需要新的劳力、资金、材料和管理者；旧有的田界、债务、人口和商路可能已经改变。一次重修碑或一项赈济记录，最能说明人们试图重新组织什么，不能证明整个地区已经复苏。

家庭的恢复尤其不均匀。有田产和亲族网络的人可能较容易重新登记或获得担保，失去成员、工具或住处的人则可能进入雇佣、流动或新的依附关系。女性、儿童和无籍者在这种过程中的经验往往很少独立留存。承认这一不均匀，比把战后写成普遍繁荣或普遍绝望更接近材料。

\subsection{遗民的选择与日常的延续}

宋遗民的诗文、墓志和不仕叙述，使失国后的道德情感得以保存。作者可能选择不仕、隐居、教学、修史或以纪念表达政治立场；这些选择值得严肃对待，却不应成为所有南方人的共同剧本。有人接受新任命，有人经营旧业，有人迁居，有人只是在不留下文字的情况下处理家庭生计。遗民文本最清楚地照亮其作者和读者群体，不能自动替代市场、乡村和港口的沉默者。

新政权同样要借助旧世界的技能。它需要熟悉税籍的吏员、掌握水道的船户、能维持市场的商人和能够安抚社区的地方中介。接收并非一张空白纸上的制度设计，而是选择、重用、限制或改名既有关系的过程。这种连续性使政权转型显得不那么突然，也使旧有不平等可能以新方式延续。

\subsection{记忆如何改变历史尺度}

崖山后来成为南宋覆亡最强的象征之一。它把海上舰队、宗室与忠义凝聚成一个可被反复讲述的场景。象征能帮助后人理解失国的重量，却也可能压缩此前数十年在江淮、四川、荆襄、闽广和沿海发生的不同危机。将崖山放回这些地区的运输、迁徙、宗教与市场网络，既不削弱其意义，也避免把整个社会的变化缩成一日。

历史记忆的价值不在于替事实盖章，而在于显示不同世代为何需要某种过去。宋遗民、元代地方社会和后来的读者，都可能以不同方式选择南宋的终结作为自己理解秩序、忠诚与地方的资源。研究者应区分这种后来的使用，和当时能由编年、文书、碑刻或考古材料确认的行动。只有这样，记忆才不会取代历史，而能成为历史的一部分。
